\section{Introduction}

Reperfusion treatment is the main acute treatment for ischaemic stroke. Intravenous thrombolysis (IVT) restores blood flow by dissolving the clot, and mechanical thrombectomy (MT) removes a clot from a large artery using an endovascular device. Both treatments are supported by randomised controlled trials, but those trials were designed to answer whether treatment works in selected patients, not to describe its effect across the full range of patients treated in routine care.

The benefit of IVT was first shown in the NINDS trial of alteplase given within 3 hours of stroke onset\cite{ninds_tissue_1995}. ECASS III extended the treatment window to 4.5 hours\cite{hacke_thrombolysis_2008}. An individual patient data meta-analysis of nine trials then showed that the benefit of alteplase falls as time to treatment increases, and that, within 4.5 hours, treatment improves the chance of a good outcome regardless of age or stroke severity\cite{emberson_effect_2014}. Beyond 4.5 hours, trials have selected patients using advanced brain imaging. In EXTEND, patients with salvageable brain tissue on perfusion imaging benefited from alteplase given 4.5 to 9 hours after onset\cite{ma_thrombolysis_2019}, a finding supported by a pooled analysis of perfusion-selected trials\cite{campbell_extending_2019}. At the mild end of the severity range, the PRISMS trial found no benefit of alteplase over aspirin in patients with minor stroke without disabling deficits, although the trial stopped early\cite{khatri_effect_2018}.

MT was established by a series of trials published in 2015. The HERMES meta-analysis of five of these trials showed that thrombectomy, mostly given with IVT, substantially reduced disability in patients with large-vessel occlusion treated mostly within 6 hours of onset\cite{goyal_endovascular_2016}. The DAWN and DEFUSE 3 trials then showed benefit up to 24 and 16 hours after onset, respectively, in patients selected by clinical--imaging or perfusion mismatch\cite{nogueira_thrombectomy_2018, albers_thrombectomy_2018}. Since then, perfusion imaging has been used increasingly to select patients for late treatment.

These trials provide strong evidence of benefit, but they leave gaps that matter for clinicians, patients and service planners. First, trials report a small number of outcomes, most often the modified Rankin Scale (mRS) at 90 days. Outcomes that matter to patients and health services during the acute admission, such as in-hospital death, length of stay and where patients go after leaving hospital, are rarely reported in detail. Second, trials are powered to detect an overall treatment effect. Subgroup results by stroke severity are therefore imprecise, and some combinations of treatment, timing and severity have little or no trial data. Third, trial populations are selected, for example by age, pre-stroke disability or imaging criteria, and may not represent all patients treated in practice.

Large national stroke registers can help fill these gaps. The Sentinel Stroke National Audit Programme (SSNAP) collects data on more than 90\% of stroke admissions in England and Wales\cite{ssnap_website}. However, simple comparisons of treated and untreated patients in registry data are prone to bias. Target trial emulation addresses this by first specifying the randomised trial that would ideally be run, and then emulating its eligibility criteria, treatment strategies, outcomes and analysis as closely as possible using observational data\cite{hernan_using_2016, hernan_target_2022}.

This study uses target trial emulation in SSNAP data to estimate the effects of IVT and MT in routine practice in more detail than existing trials allow. We emulated a separate trial for each combination of stroke severity (NIHSS 1--5, 6--10, 11--15, 16--20 and 21--42), treatment (IVT alone, MT alone, or both), treatment timing (within or after 4.5 hours of onset) and, for late treatment, use of perfusion imaging. For each trial we report five outcomes at discharge from hospital: good functional outcome, in-hospital death, disability (mRS), length of stay and discharge destination. We address three questions: how the association between treatment and outcome varies with stroke severity; how it varies with the type and timing of treatment; and whether the pattern is consistent across outcomes.
